\section{LLM, RAG und Tool Use (Kap. 03)}
\label{sec:llm-rag-tool-use}

\subsection{LLM- und RAG-Entscheidung}

Als Standardmodell wird \texttt{gpt-5.6-luna} verwendet, alternativ kann ein OpenAI-kompatibler HAW-Endpunkt
verwendet werden. Für den Lernchat ist diese Wahl plausibel als Kosten- und
Latenzkompromiss, weil im Normalfall drei Spezialistenaufrufe und bei der
voreingestellten zweiten Koordinationsrunde sechs entstehen. Konfigurierbar sind
bis zu zwölf Aufrufe, jeweils zuzüglich der internen Toolrunden. Benötigt werden
JSON-Modus und Function Calling, nicht
maximale Modellgröße. Modellabhängige Kompatibilität ist sichtbar: Reasoning-Modelle erhalten keinen
Temperaturparameter. Für GPT-5-Function-Calls setzt der Adapter im
Chat-Completions-Endpunkt \texttt{reasoning\_effort=none}. Ein empirischer Qualitätsvergleich verschiedener Modelle
fehlt. Ihre Eignung ist daher nicht empirisch belegt.

RAG ist fachlich begründet, weil private Lernmaterialien weder zuverlässig
im Modellwissen liegen noch vollständig in jeden Prompt passen. Lewis et al.
beschreiben die Verbindung parametrischen Modellwissens mit explizitem,
nichtparametrischem Speicher \parencite{lewisEtAl2020rag}. StudyMummy indexiert
in einer Offline-Phase Dokumentchunks mit dem konfigurierten Embedding-Modell
(Default \texttt{text-embedding-3-small}) in \texttt{pgvector}. Online folgen
Retrieval, Prompt-Augmentation und Antwortgenerierung: Pro Frage werden die drei
ähnlichsten, nutzerisolierten Chunks geladen. Der Nutzen ist
projektspezifische Grundierung, keine Wahrheitsgarantie. Grenzen sind starres
Zeichen-Chunking, fehlende Ähnlichkeitsschwelle und Quellenangabe, kein
Reranking und stilles Degradieren bei Fehlern. Auch korrekter Kontext schließt
Halluzinationen nicht aus. Embedding und Retrieval erhöhen zudem Latenz und
Kosten.

\subsection{Tools und Agentic Pattern}

Die Registry enthält sechs Werkzeuge. \Cref{lst:tools} zeigt die für Function
Calling relevanten, vereinfachten Schnittstellen. Vier davon liegen abhängig
von der geplanten Aktion im effektiven Chat-Scope.

\begin{lstlisting}[language=Python,float=htbp,caption={Registrierte Werkzeugsignaturen und effektiver Chat-Scope.},label={lst:tools}]
# vom Chat-Workflow freigebbar
evaluate_answer(task_id: str, user_answer: str,
                 expected_concept: str, help_level: int = 1) -> dict
update_learning_profile(tag: str, score: float,
                        error_pattern: str = "") -> dict
award_coins_and_exp(amount: int, reason: str,
                    task_id: str | None = None) -> dict
add_calendar_note(title: str, content: str,
                  start_time: str, end_time: str) -> dict

# registriert, aber durch SAFE_TOOL_POLICY im Chat nicht erreichbar
generate_quiz_questions(topic: str, num_questions: int = 5,
                         difficulty: int = 2) -> dict
create_cheatsheet(user_id: str, session_id: str,
                  topics: list[str]) -> dict
\end{lstlisting}

Das äußere Muster ist eine begrenzte
\emph{Plan--Act--Observe--Review--Replan}-Schleife. Der Tutor alterniert intern
Function Calls und Beobachtungen. Der Reviewer kann den Entwurf akzeptieren,
toolfrei revidieren lassen oder Beobachtung und Kritik an den Planner
zurückgeben. Dies überträgt ReActs Kopplung von Reasoning und Handeln auf den
typisierten MAS-Kontrollfluss \parencite{yaoEtAl2023react}, ohne interne
Gedankengänge zu veröffentlichen. \texttt{tool\_calls} enthält nur erfolgreiche
Funktionen. Alle Versuche werden separat als \texttt{ToolObservation} mit
Status und gekürzter Resultatvorschau an Reviewer und API gemeldet. Plan,
sanitisierte Nachrichtenmetadaten, lokale Agentenzustände und
Koordinationsrunden sind beobachtbar, werden aber weiterhin nicht als
vollständige Episode persistiert.
